\chapter{Kolmogorov-Chentsov Theorem}
\label{chap:kolmogorov_chentsov}

We follow the proof of the Kolmogorov-Chentsov theorem from \cite{kratschmer2023kolmogorov}.
That proof notably uses the chaining technique developed by Talagrand \cite{talagrand2022upper}.

That theorem is about stochastic processes $X : T \to \Omega \to E$, where $\Omega$ is a measurable space with a probability measure $\mathbb{P}$, the index set $T$ is a metric space with distance $d_T$, and $E$ is also a metric space with distance $d_E$, on which we put the Borel $\sigma$-algebra.

The main result is Theorem~\ref{thm:countable_set_bound}.
Under an assumption on the covering number of $T$, for a process $X$ that satisfies the Kolmogorov condition $\mathbb{E}[d_E(X_s, X_t)^p] \le M d_T(s, t)^q$ (see Definition~\ref{def:IsKolmogorovProcess}),the theorem gives a finite bound on the expectation of the supremum of the ratio
\begin{align*}
  \mathbb{E}\left[ \sup_{s, t \in T'} \frac{d_E(X_s, X_t)^p}{d_T(s, t)^{\beta p}} \right]
  \: ,
\end{align*}
for $T'$ a countable subset of $T$.
As a corollary, we obtain that there exists a modification of $X$ with Hölder continuous paths.

In Lean, we will use the typeclass \texttt{PseudoEMetricSpace} for both $T$ and $E$ as long as possible, and then specialize to \texttt{EMetricSpace} (or perhaps even \texttt{MetricSpace}) when we need the stronger properties of a metric space.
For example, to prove the existence of a modification of a stochastic process, we will eventually use the fact that $d_E(x, y) = 0$ implies $x = y$, which does not hold in a pseudo-metric space.
All distances will be expressed with \texttt{edist}, which takes values in \texttt{ENNReal}, and the integrals refer to Lebesgue integrals.

\section{Covers and covering numbers}

Let $(E, d_E)$ be a pseudo-metric space.

\inputleannode{def:IsCover}


\inputleannode{def:externalCoveringNumber}


\inputleannode{def:internalCoveringNumber}


\inputleannode{def:IsSeparated}


\inputleannode{def:packingNumber}


\inputleannode{lem:externalCoveringNumber_le_internalCoveringNumber}




\inputleannode{lem:internalCoveringNumber_le_packingNumber}




\inputleannode{lem:packingNumber_two_le_externalCoveringNumber}




\inputleannode{lem:internalCoveringNumber_eq_one_of_diam_le}




\inputleannode{lem:externalCoveringNumber_mono}




\inputleannode{lem:internalCoveringNumber_subset_le}




\subsection{Covering number and volume}

In this section $E$ is a finite dimensional inner product space, with dimension $d$.

\inputleannode{lem:volume_le_of_isCover}



The volume of $B_\varepsilon$ is given by the following formula (see \href{https://leanprover-community.github.io/mathlib4_docs/Mathlib/MeasureTheory/Measure/Lebesgue/VolumeOfBalls.html#InnerProductSpace.volume_closedBall}{InnerProductSpace.volume\_closedBall}).
\begin{align*}
  V(B_\varepsilon) = \frac{\pi^{d/2}}{\Gamma(d/2 + 1)} \varepsilon^d
  \: .
\end{align*}


\inputleannode{lem:volume_le_externalCoveringNumber_mul}




\inputleannode{lem:le_volume_of_isSeparated}




\inputleannode{lem:packingNumber_mul_le_volume}




\inputleannode{lem:volume_div_le_internalCoveringNumber}




\inputleannode{lem:internalCoveringNumber_le_volume_div}




\inputleannode{lem:internalCoveringNumber_closedBall_ge}




\inputleannode{lem:internalCoveringNumber_closedBall_le}




\subsection{Bounded internal covering number}

\inputleannode{def:HasBoundedInternalCoveringNumber}


\inputleannode{lem:hasBoundedInternalCoveringNumber_unitInterval}




\inputleannode{lem:hasBoundedInternalCoveringNumber_subset}




\section{Chaining}

\subsection{Chaining sequence}


\inputleannode{def:nearestPt}


\inputleannode{lem:dist_nearestPt_le}




\inputleannode{lem:dist_nearestPt_of_isCover}




\inputleannode{def:chainingSequence}


\inputleannode{lem:chainingSequence_mem}




\inputleannode{lem:dist_chainingSequence_add_one}




\inputleannode{lem:dist_chainingSequence_le_sum}




\inputleannode{lem:dist_chainingSequence_le}




\inputleannode{cor:dist_chainingSequence_pow_two_le}




\subsection{A subset of pairs}

We will be interested in bounding expressions of the form $\sup_{s,t\in J, d_T(s,t) \le c} d_E(f(s), f(t))$ for a finite set $J$ and some function $f : T \to E$.
This is a supremum over pairs in $J$ and there could be $\vert J \vert^2$ such pairs.
We will build a subset $K$ of $J^2$ which is much smaller, of size linear in $\vert J \vert$, such that its points are not too far apart and
\begin{align*}
  \sup_{s,t\in J, d_T(s,t) \le c} d_E(f(s), f(t))
  & \le 2 \sup_{(s,t) \in K} d_E(f(s), f(t))
  \: .
\end{align*}
The pairs $(s, t) \in K$ will still be close together, in the sense that $d_T(s, t) \le c n$ for some $n$ that is logarithmic in the size of $J$.

For $t \in V \subseteq T$ and $u\ge 0$, we denote by $B_V(t, u)$ the closed ball with center $t$ and radius $u$ in $V$.
That is, $B_V(t, u) = \{s \in V \mid d_T(s, t) \le u\}$.

We want to cover $J$ with balls that have radius logarithmic in the number of points of the ball.

\inputleannode{def:logSizeRadius}


\inputleannode{lem:card_logSizeRadius_ge}




\inputleannode{lem:card_logSizeRadius_le}




\inputleannode{def:logSizeBallSequence}

A log-size ball sequence gives a partition of $J$ into sets which are contained in balls of radius $(r_i - 1)c$ around $t_i$, and satisfy cardinality constraints.


\inputleannode{lem:logSizeRadius_logSizeBallSequence_le}




\inputleannode{lem:logSizeBallSequence_V_anti}




\inputleannode{lem:logSizeBallSequence_eq_zero}




\inputleannode{lem:logSizeBallSequence_disjoint_B}




\inputleannode{def:pairSet}


\inputleannode{lem:card_pairSet_le}




\inputleannode{lem:dist_le_of_mem_pairSet}




\inputleannode{lem:sup_dist_le_two_mul_sup_dist_pairSet}




\inputleannode{lem:pair_reduction}







\section{Chaining for stochastic processes under Kolmogorov conditions}

\subsection{Kolmogorov condition}

\inputleannode{def:IsKolmogorovProcess}

Remark: the measurability condition on the pair would be implied by the measurability of $X_t$ for all $t \in T$ if we assumed that $E$ is separable (\texttt{SecondCountableTopology} in Lean), which implies that the Borel $\sigma$-algebra on the product is equal to the product of the Borel $\sigma$-algebras.
We follow \cite{kratschmer2023kolmogorov} and do not require separability.


\inputleannode{lem:IsKolmogorovProcess.edist_eq_zero}




\inputleannode{lem:IsKolmogorovProcess.lintegral_sup_rpow_edist_eq_zero}




\paragraph{Measurability}

\inputleannode{lem:IsKolmogorovProcess.aemeasurable}




\inputleannode{lem:aemeasurable_pair_of_aemeasurable}




\inputleannode{lem:IsKolmogorovProcess.aemeasurable_edist}



\paragraph{Distance bounds}

\inputleannode{lem:integral_sup_rpow_dist_le_card_mul_rpow}




\inputleannode{lem:integral_sup_rpow_dist_of_dist_le}




\subsection{Bound for a set of points that are close together}

For a finite index set $T$, we want to obtain a bound on
\begin{align*}
  \mathbb{E}\left[ \sup_{s, t \in T; d_T(s, t) \le \delta} d_E(X_s, X_t)^p \right] \: .
\end{align*}
Note the condition that the supremum is taken over pairs $(s, t)$ such that $d_T(s, t) \le \delta$.

We consider covers of $T$ at different scales. $C_n$ is a finite $\varepsilon_n$-cover of $T$ with $\varepsilon_n = \varepsilon_0 2^{-n}$.
$T$ is equal to $C_k$ for some $k$ large enough, so the supremum over $T$ is a supremum at that scale $k$.
We will change scale to some $m \le k$ that depends on the distance bound $\delta$ ($m$ is of order $\log_2 \delta$) and consider the supremum over $C_m$ (plus a term due to the scale change).
Then for the supremum over a set in $C_m$, we use the reduction in the number of pairs of Lemma~\ref{lem:pair_reduction}.

\inputleannode{lem:scale_change}




\inputleannode{cor:scale_change_rpow}





\subsubsection{First term}


\inputleannode{lem:integral_sup_rpow_dist_cover_of_dist_le}



\inputleannode{lem:integral_sup_rpow_dist_cover_rescale}





\subsubsection{Second term}


\inputleannode{lem:integral_sup_rpow_dist_succ}





\paragraph{Case $p \ge 1$}


\inputleannode{lem:integral_sup_dist_le_sum_rpow}




\inputleannode{lem:integral_sup_rpow_dist_le_sum}




\inputleannode{lem:integral_sup_rpow_dist_le_of_minimal_cover}




\inputleannode{cor:integral_sup_rpow_dist_le_of_minimal_cover_two}





\paragraph{Case $p \le 1$}


\inputleannode{lem:integral_sup_dist_le_sum_rpow_of_le_one}




\inputleannode{lem:integral_sup_rpow_dist_le_sum_of_le_one}




\inputleannode{lem:integral_sup_rpow_dist_le_of_minimal_cover_of_le_one}




\inputleannode{cor:integral_sup_rpow_dist_le_of_minimal_cover_two_of_le_one}




\paragraph{Any $p>0$}


\inputleannode{def:Cp}


\inputleannode{lem:second_term_bound}





\subsubsection{Putting it all together}


\inputleannode{lem:lintegral_sup_cover_eq_of_lt_iInf_dist}




\inputleannode{thm:finite_set_bound_of_dist_le_of_diam_le}




\inputleannode{thm:finite_set_bound_of_dist_le_of_le_diam}




\inputleannode{cor:finite_set_bound_of_dist_le_of_le_diam_bis}




\inputleannode{cor:finite_set_bound_of_dist_le}







\section{Kolmogorov-Chentsov Theorem}


\subsection{Sets with bounded internal covering number}

TODO: change the proofs here to avoid $s \ne t$ and use instead properties of processes satisfying the Kolmogorov condition for exponents $(p,q)$.

\inputleannode{lem:integral_div_dist_le_sum_integral_dist_le}




\inputleannode{def:L}


\inputleannode{lem:L_lt_top}




\inputleannode{lem:finite_set_bound}




\inputleannode{thm:countable_set_bound}




\begin{corollary}\label{cor:countable_set_bound_of_le}
Under the same assumptions as in Theorem~\ref{thm:countable_set_bound}, for every countable subset $T' \subseteq T$ with positive diameter, for $L(T, c_1, d, p, q, \beta) < \infty$ the same constant,
\begin{align*}
  \mathbb{E}\left[ \sup_{s, t \in T';\: d_T(s, t) \le \delta} d_E(X_s, X_t)^p \right]
  \le M L(T, c_1, d, p, q, \beta) \delta^{\beta p}
  \: .
\end{align*}
\end{corollary}

\begin{proof}
  \uses{thm:countable_set_bound}
Immediately follows from Theorem~\ref{thm:countable_set_bound}.
\end{proof}


\inputleannode{lem:holder_modification_single}




\inputleannode{thm:holder_modification}





\subsection{Localized Kolmogorov-Chentsov theorem}

\inputleannode{def:HasBoundedCoveringNumberCover}


\inputleannode{lem:hasBoundedCoveringNumberCover_nnreal}

Note: in the Lean code, we proved this result for weaker constants: $c_n = 3n$.



We say that a function is \emph{locally} Hölder continuous of order $\gamma$ if for any point $x$ in its domain there is a neighborhood of $x$ on which the function is Hölder continuous of order $\gamma$.

\inputleannode{thm:localized_holder_modification}




\inputleannode{thm:localized_holder_modification_sup}


